\def\stat{dulin}





\def\tit{ВЫБОР ТЕХНОЛОГИЧЕСКИХ РЕШЕНИЙ ДЛЯ~ПОДДЕРЖКИ ПРОЦЕССА СИНТЕЗА 


ГЕОДАННЫХ ИНФРАСТРУКТУРЫ ЖЕЛЕЗНОДОРОЖНОГО ТРАНСПОРТА}





\def\titkol{Выбор технологических решений для~поддержки процесса синтеза 


геоданных} % инфраструктуры железнодорожного транспорта}





\def\aut{И.\,Н.~Розенберг$^1$, С.\,К.~Дулин$^2$}





\def\autkol{И.\,Н.~Розенберг, С.\,К.~Дулин}





\titel{\tit}{\aut}{\autkol}{\titkol}





\index{Розенберг И.\,Н.}


\index{Дулин С.\,К.}


\index{Rozenberg I.\,N.}


\index{Dulin S.\,K.}





%{\renewcommand{\thefootnote}{\fnsymbol{footnote}}


%\footnotetext[1]


%{Исследование выполнено с~использованием ЦКП <<Информатика>> ФИЦ ИУ РАН.}} 





\renewcommand{\thefootnote}{\arabic{footnote}}


\footnotetext[1]{Научно-исследовательский и~проектно-конструкторский институт информатизации, автоматизации и~связи 


на железнодорожном транспорте (АО НИИАС), \mbox{I.Rozenberg@vniias.ru}}


\footnotetext[2]{Федеральный исследовательский центр <<Информатика и~управление>> Российской 


академии наук; Научно-исследовательский и~проектно-конструкторский институт 


информатизации, автоматизации и~связи на железнодорожном транспорте (АО НИИАС), 


\mbox{skdulin@mail.ru}}





 \label{st\stat}





 \thispagestyle{headings}


 


 \vspace*{-14pt}


  





  


  


  \Abst{Слияние, или синтез, данных (data fusion)~--- процесс, при котором для 


анализа используются данные, объединенные из нескольких источников для получения 


более со\-гла\-су\-ющей\-ся, точ\-ной и~полезной информации. Получение и~слияние данных от 


датчиков определяется как мультисенсорное слияние данных. Самым ярким примером 


слияния данных служит собственно человек, который постоянно комбинирует данные, 


по\-сту\-па\-ющие от всех органов чувств для выполнения необходимых действий 


в~окру\-жа\-ющей среде. В~част\-ности, человек не смог бы водить машину без слияния 


зрительных, слуховых данных и~данных моторики. В~геоинформатике вместо синтеза 


геоданных час\-то используется термин интеграция геоданных. Современные географические информационные системы 


могут 


комбинировать различные наборы геоданных так, что полученные в~результате наборы 


геоданных содержат атрибуты и~метаданные, которых нет у~точек в~ком\-би\-ни\-ру\-емых 


наборах данных.


  В~\mbox{статье} анализируются категории и~модели синтеза геоданных, ис\-поль\-зу\-емые для 


обеспечения поддержки актуальности сис\-те\-мы пространственных данных 


в~инфраструктуре железнодорожного транс\-порта.}


  


  \KW{синтез данных; категории синтеза данных; модель синтеза данных}


  


\DOI{10.14357\!/\!08696527230205}  





\vskip 10pt plus 9pt minus 6pt





 \thispagestyle{headings}


 


 \vspace*{-18pt}


      


  \section{Введение}


  


  \vspace*{-3pt}


  


  Тема синтеза геоданных, или интеграции геоданных, объединенных из 


нескольких источников, не нова, но современные технологии интеграции 


геоданных предостав\-ля\-ют возможности для улучшения этого процесса~[1]. 


Интеграция геоданных в~распределенных информационных средах 


с~функциональной со\-вмес\-ти\-мостью, основанной на открытых стандартах, 


радикально меняет классические области процесса интеграции, предлагая 


совершенно новые способы распознавания отношений. Геопространственные 


мероприятия, основанные на местоположении и~времени, в~на\-сто\-ящее время 


стали наиболее распространенными, хотя ге\-те\-ро\-ген\-ность геоданных сильно 


услож\-ня\-ет проведение интеграции. 


  


  Открытый геопространственный консорциум (Open Geospatial Consortium, OGC) стал международной 


некоммерческой добровольной организацией по стандартам отраслевого 


консенсуса, которая предостав\-ля\-ет платформу для обсуж\-де\-ния и~предлагает 


верифицированные процессы для совместной разработки бес\-плат\-ных 


и~общедоступных спецификаций интерфейса (открытых стандартов). Эти 


открытые стандарты обеспечивают более легкий до\-ступ 


к~геопространственной информации и~ее использование, интеграцию, 


а~так\-же улучшенную функциональную совместимость геопространственных 


технологий (на любом устройстве, платформе, сис\-те\-ме, сети или 


предприятии) для удовлетворения по\-треб\-но\-стей мирового сообщества. 


Открытые стандарты OGC широко внед\-ре\-ны на рынке геоинформационных 


услуг и~помогают развивать распределенный и~компонентный синтез 


решений, которые поддерживают гео-веб, беспроводные и~геолокационные 


сервисы, а~также более широкие правительственные и~деловые 


 информационно-тех\-но\-ло\-ги\-че\-ские пред\-при\-ятия по всему миру.


  


     Для выполнения своей деятельности консорциум OGC проводит три 


вида программ.


     \begin{enumerate}[1.]


\item Программа спецификаций OGC способствует созданию официальных 


комитетов, рабочих групп и~групп с~особыми интересами, основанных на 


консенсусе, которые организуют форум для отраслевых предприятий, 


исследовательских и~пользовательских сообществ, обеспечивая совместное 


выявление, приоритизацию и~продвижение решений с~\mbox{целью} удовле\-тво\-ре\-ния 


по\-треб\-но\-стей в~стандартах. 


\item Программа функциональной совместимости и~синтеза решений 


способствует быст\-ро\-му прототипированию, тестированию и~валидации 


новых стандартов с~по\-мощью быст\-ро\-раз\-ви\-ва\-ющих\-ся испытательных 


стендов, экспериментов, пилотных инициатив и~связанных с~ними  


тех\-ни\-ко-эко\-но\-ми\-че\-ских обоснований. 


\item Программы OGC по ин\-фор\-ма\-ци\-он\-но-разъ\-яс\-ни\-тель\-ной работе 


и~адап\-та\-ции в~обществе~--- это программы повышения осве\-дом\-лен\-ности 


и~внед\-ре\-ния стандартов OGC в~геоинформационном сообществе. 


\end{enumerate}





  Следует отметить, что без открытых стандартов геоданные и~сервисы 


слож\-но синтезировать и~масштабировать. Данные, приложения и~сервисы, 


соответствующие стандартам, создают автоматизированную 


и~интероперабельную среду синтеза геоданных, обеспечивая совместный 


доступ к~геоданным и~сервисам их обработки. 


  


  Стандартизованная открытая структура синтеза геоданных широко 


распространена в~среде потребителей георесурсов. Так, стандарт обработки 


карты в~сети (Web Map Service~[2]~--- WMS) поддерживает синтез карт. 


Стандарт WMS синтезирует карты в~виде уровней изображений (рис.~1).














  Далее в~работе анализируются категории и~модели синтеза геоданных, 


ис\-поль\-зу\-емые для обеспечения поддержки ак\-ту\-аль\-ности сис\-те\-мы 


пространственных данных инфраструктуры железнодорожного транспорта.





  


  \section{Категории синтеза данных}


  


  \begin{figure} %fig1


    \vspace*{1pt}


  \begin{center}


 \mbox{%


 \epsfxsize=128mm 


\epsfbox{dul-1.eps}


 }


\end{center}


\vspace*{-11pt}


\Caption{Вариант геоинформационного синтеза, поддерживающего WMS-стан\-дарт OGC. 


Геоданные последовательно синтезируются в~единственную визуализацию}


\end{figure}


  


  Синтез данных, как определено OGC, выглядит сле\-ду\-ющим образом:


   


   <<\textit{Синтез данных (data fusion)~--- это акт или процесс 


объединения или связывания данных или информации, ка\-са\-ющих\-ся одной 


или нескольких сущностей, рассматриваемых в~явной или неявной 


структуре знаний, для улучшения возможностей обнаружения, 


идентификации или характеристик ис\-сле\-ду\-емых сущностей}>>~[3].


     


     При обсуждении синтеза данных обычно используются три категории 


(рис.~2):


     \begin{enumerate}[(1)]


\item синтез наблюдений (показаний датчиков); 


\item синтез объектов; 


\item синтез решений.


\end{enumerate}





  Фундаментальным понятием в~любом типе синтеза данных остается 


построение ассоциаций между одним или несколькими элементами данных. 


Во многих случаях эта ассоциация может так\-же включать дополнительную 


обработку и~абстракцию и~приводить к~элементу данных нового типа. Это 


понятие приводит к~категориям синтеза данных, основанным на типе 


связанных элементов данных: синтез наблюдений, синтез объектов и~синтез 


решений (см.\ рис.~2).


  


  Обычно термин <<синтез датчиков>> используется для обозначения 


синтеза наблюдений. Синтез наблюдений~--- более точный 


и~предпочтительный термин. Синтез объектов представляется более 


проб\-ле\-ма\-тич\-ным, чем синтез наблюдений. Эти два термина имеют раз\-ную 


семантику в~отдельных областях обработки изображений и~ГИС. 


  


      


      \begin{figure} %fig2


       \vspace*{1pt}


  \begin{center}


 \mbox{%


 \epsfxsize=61.801mm 


\epsfbox{dul-2.eps}


 }


\end{center}


\vspace*{-11pt}


      \Caption{Архитектура сервисов синтеза}


      \end{figure}


      


  Определения трех категорий синтеза, используемых в~настоящей работе:


  \begin{enumerate}[(1)]


\item синтез наблюдений (показаний датчиков~[4])~--- это объединение 


нескольких сенсорных измерений одних и~тех же явлений или событий 


в~комбинированное наблюдение и~анализ сигнатуры измерений, а~также 


синтез варь\-и\-ру\-ющих\-ся сенсорных измерений различных наблюдаемых 


свойств и~достаточно пол\-но описанных наблюдений, вклю\-чая 


неопределенности;


\item синтез объектов~--- это обработка наблюдений для получения 


семантических признаков более высокого порядка и~обработка этих 


признаков. Синтез объектов улучшает понимание оперативной ситуации 


и~оценку потенциальных угроз и~воздействий для выявления, 


классификации, ассоциирования и~агрегирования объектов. Процессы 


синтеза объектов и~при\-зна\-ков включают так\-же обобщение и~объединение 


при\-зна\-ков; 


\item синтез решений~--- это акт или процесс поддержки спо\-соб\-ности 


человека принимать решение путем предостав\-ле\-ния среды совместимых 


сетевых услуг для оценки ситуации, оценки воздействия и~поддержки 


принятия решений с~использованием информации от нескольких датчиков 


и~другой обработанной информации.


\end{enumerate}





  На рис.~3 представлена функциональная схема узла синтеза решений 


с~пояснением действий в~таб\-ли\-це~[3].


  





     


     \section{Модель синтеза данных}


     


  В качестве примера реализации успеш\-но\-го синтеза решений мож\-но 


привести работу <<Межведомственной сис\-те\-мы ситуационной 


осве\-дом\-лен\-ности>> (Multi-Agency Situational Awareness System, МАSАS) канадских агентств по управ\-ле\-нию 


чрезвычайными ситуациями~[5] по созданию национального потенциала для 


обмена информацией о~геопространственных чрезвычайных инцидентах. 


Основанная на открытых и~совместимых стандартах сис\-те\-ма объединяет 


информацию из нескольких источников для под\-держ\-ки процесса принятия 


решений, улучшая ситуационную осведомленность участ\-ву\-ющих 


организаций. Архитектура \mbox{МАSAS} использовалась для разработки 


и~развертывания компонентных сис\-тем \mbox{МАSAS} для обмена 


информацией об инцидентах. Существует возможность расширения 


\mbox{МАSAS} на ряд других приложений для объединения решений, таких 


как синтез информации о~реагировании на чрезвычайные ситуации 


и~трансграничный синтез информации.





\begin{figure} %fig3


 \vspace*{1pt}


  \begin{center}


 \mbox{%


 \epsfxsize=125.972mm 


\epsfbox{dul-3.eps}


 }


\end{center}


\vspace*{-11pt}


\Caption{Функциональная схема узла синтеза решений}


\end{figure}








  


  Другим примером функционирующего синтеза решений служит 


интегрированная сис\-те\-ма оповещения IPAWS (Integrated Public Alerting and Warning System), предназначенная для 


повышения общественной безопас\-ности по\-сред\-ст\-вом распространения 


экстренных сообщений в~социальной среде по всем воз\-мож\-ным устройствам 


связи~[6]. В~IPAWS в~рамках нескольких проектных инициатив 


используются специальные технологии обмена сообщениями для улуч\-ше\-ния 


и~расширения возможностей оповещения ава\-рий\-но-спа\-са\-тель\-ны\-ми 


службами, менеджерами по инцидентам и~государственными должностными 


лицами на всех уровнях государственного управ\-ле\-ния. 





\begin{table}\small %tabl


\begin{center}


\begin{tabular}{|p{41mm}|p{80mm}|}


\multicolumn{2}{c}{Информационные потоки узлов синтеза решений на рис.~3}\\


\multicolumn{2}{c}{\ }\\[-6pt]


\hline


\multicolumn{1}{|c|}{\textbf{Информационный поток}}&\multicolumn{1}{c|}{\textbf{Описание 


информационного потока}}\\


\hline


\textit{Команды} из командных узлов&Структурированная информация, на основе 


которой должен действовать узел синтеза решений\\


\hline


\textit{Сообщения} командным узлам&Информация о структуре, полученная в~результате 


синтеза на узле синтеза решений \\


\hline


\textit{Открытый исходный} код расширенного поиска&Сбор неструктурированной 


информации с~помощью узла синтеза решений из интернета (или через узел сбора данных с~открытым исходным кодом) решений \\


\hline


\textit{Данные} от датчика и~синтез показаний&Структурированная информация, 


полученная в~результате синтеза наблюдений датчиков и~анализа объектов решений \\


\hline 


\textit{Синтез заданий} датчику и~его показаний&Структурированные сообщения, 


за\-пра\-ши\-ва\-ющие результаты наблюдений, включая запрос на генерацию синтеза решений 


\\


\hline


\textit{Команды} диспетчеризации операций &Структурированная информация, 


на\-прав\-ля\-ющая де\-я\-тель\-ность диспетчеризации операций решений \\


\hline


\textit{Сообщения} от диспетчеризации операций &Структурированная информационная 


отчетность о~со\-сто\-янии диспетчерской и~оперативной де\-я\-тель\-ности решений \\


\hline


\textit{Расширенный поиск} на дру\-гих операционных узлах &Поисковые запросы 


к~операционному узлу, распознаваемому узлом синтеза решений решений \\


\hline


\textit{Сообщения} на другие операционные узлы&Структурированная информация, 


пе\-ре\-да\-ва\-емая на операционный узел, распознаваемый узлом синтеза решений решений \\


\hline


\textit{Структурированные данные} из других узлов операций&Структурированная 


информация, полученная от операционных узлов, распознаваемых узлом синтеза решений 


решений \\


\hline


\end{tabular}


\end{center}


\vspace*{-9pt}


\end{table}


  


  Одной из наиболее значимых моделей синтеза считается модель JDL~[7]. 


На рис.~4 изображен информационный поток в~модели, объединяющей 


отмеченные на рис.~1 три категории синтеза данных:


  \begin{itemize}


\item \textbf{Синтез наблюдений} (уровень~0).


\item \textbf{Синтез объектов} (уровень~1).


\item \textbf{Оценка ситуации} (уровень~2): восприятие 


элементов окру\-жа\-ющей среды в~пределах объема времени и~пространства, 


понимание их значения и~проекция их статуса на ближайшее будущее. 


\item \textbf{Оценка воздействия} (уровень~3): процесс 


рассмотрения последствий пред\-ла\-га\-емых действий для людей и~их среды 


при воз\-мож\-ности их корректировки. Он применяется на всех уровнях 


принятия решений~--- от планирования до конкретных проектов. 


\item \textbf{Совершенствование} процессов (уровень~4): 


адап\-тив\-ный сбор и~обработка данных для под\-держ\-ки целей процесса. 


\end{itemize}





  Осознание и~оценка ситуации~--- это комбинированный продукт 


восприятия, понимания и~проекции, по\-лу\-ча\-ющий\-ся в~результате обработки 


про\-стран\-ст\-вен\-ных, временн\'{ы}х и~семантических атрибутов метаданных 


объектов.


  


  Современные концептуальные модели синтеза информации, в~том чис\-ле 


и~модель JDL, не учитывают тот факт, что источники информации час\-то 


базируются на раз\-ных онтологических основаниях~[7], поэтому считается, 


что модель JDL, которая обслуживает общие ссылки на пространство 


и~время, должна быть дополнена условным аспектом общего 


онтологического вы\-рав\-ни\-ва\-ния. Онтологическое вы\-рав\-ни\-ва\-ние~--- это акт 


уста\-нов\-ле\-ния отношения соответствия между двумя или более символами из 


раз\-ных онтологий для тех символов понятий, которые семантически 


иден\-тич\-ны или похожи.


      


 


\begin{figure} %fig4


 \vspace*{1pt}


  \begin{center}


 \mbox{%


 \epsfxsize=77.881mm 


\epsfbox{dul-4.eps}


 }


\end{center}


\vspace*{-11pt}


\Caption{Модель синтеза данных JDL}


\end{figure}





  Информация, используемая в~моделях синтеза данных, может 


рас\-смат\-ри\-вать\-ся в~двух широких категориях: структурированная 


и~неструктурированная информация~[8]. При этом структурированная 


информация определяется как со\-от\-вет\-ст\-ву\-ющая структурам уста\-нов\-лен\-ных 


стандартов. Существующие стандарты OGC для структурированной 


информации вклю\-ча\-ют в~себя спецификацию OGC признаков, покрытий, 


наблюдений и~гео\-мет\-рии.


  





\begin{figure} %fig5


 \vspace*{1pt}


  \begin{center}


 \mbox{%


 \epsfxsize=128mm 


\epsfbox{dul-5.eps}


 }


\end{center}


\vspace*{-12pt}


\Caption{Синтез решений на основе коммутации геоданных из разных источников}


\vspace*{-6pt}


\end{figure}


     


  На рис.~5 показан пример синтеза решений~[9], который связан 


  с~анализом тенденций и~причин паводка в~районе слияния рек на основе 


геоданных из нескольких источников. Синтез решений, проводимый в~этом 


исследовании, пред\-став\-ля\-ет собой комплексную операцию, вклю\-ча\-ющую как 


геоданные, по\-лу\-ча\-емые от человека с~мобильных устройств, так и~геоданные 


общей обстановки из цент\-ра МЧС. Ис\-поль\-зу\-ющи\-еся геоданные 


пред\-став\-ля\-ют собой как структурированную, так и~неструктурированную 


информацию~[10]. Таким образом, информация, доступная для 


осуществления синтеза решений,~--- это коллекция сведений 


интеллектуального корпоративного капитала, поскольку источниками таких 


сведений вы\-сту\-па\-ют люди, документы, оборудование или технические 


датчики. Эти сведения могут быть сгруппированы сле\-ду\-ющим образом:


  \begin{itemize}


\item сведения, получаемые от людей;


\item геопространственные сведения, накопленные в~виде георесурсов;


\item сведения в~виде сигналов датчиков; 


\item сведения в~виде данных измерения разнообразных параметров;


\item сведения о геоданных открытого доступа произвольного формата;


\item сведения о на\-уч\-но-тех\-ни\-че\-ских разведывательных геоданных. 


\end{itemize}





  Развитие синтеза решений требует разработки стандартов, таких как 


методы отображения схемы с~соответствующей идентификацией правил 


отображения, и~увеличения ак\-цен\-та на обрабатывание ассоциаций, поскольку 


идентификация ассоциации между объектами лежит в~основе синтеза. Самая 


эффективная среда для реализации указанных категорий синтеза~--- это 


архитектура с~распределенными базами данных и~сервисами, основанными 


на об\-щем яд\-ре форматов геоданных, алгоритмов и~приложений~[11]. 


  


\section{Интеграция геоданных в~комплексной системе 


пространственных данных инфраструктуры железнодорожного 


транспорта}





  Существующие подходы к~реализации синтеза решений найдут свое 


применение в~информационном обес\-пе\-че\-нии проектирования, строительства, 


реконструкции и~капитального ремонта объектов железнодорожной 


инфраструктуры. Комплексная сис\-те\-ма пространственных данных 


инфраструктуры железнодорожного транс\-пор\-та (КСПД ИЖТ) 


разрабатывается в~соответствии с~распоряжением президента ОАО <<РЖД>> 


от 3~декабря 2010~года №\,2511р <<О~создании комплексной сис\-те\-мы 


пространственных данных инфраструктуры железнодорожного 


транс\-пор\-та>>. Гетерогенность геоданных, подлежащих комплексированию 


и~синтезу в~КСПД ИЖТ, выдвигает на передний план решение задачи 


гео\-ин\-тер\-опе\-ра\-бель\-ности отраслевого масштаба~[12, 13] как задачи 


совместного согласованного использования геоданных, полученных из 


разных информационных источников (рис.~6).


     


\begin{figure} %fig6


 \vspace*{1pt}


  \begin{center}


 \mbox{%


 \epsfxsize=127.68mm 


\epsfbox{dul-6.eps}


 }


\end{center}


\vspace*{-11pt}


\Caption{Функционал КСПД ИЖТ~[7]}


\end{figure}


  


  КСПД ИЖТ~--- это ин\-фор\-ма\-ци\-он\-но-тех\-но\-ло\-ги\-че\-ская сис\-те\-ма 


централизованного сбора, интеграции, хранения и~анализа  


ко\-ор\-ди\-нат\-но-при\-вя\-зан\-ных сведений об объектах инфраструктуры 


железнодорожного транспорта. Разрабатываемая сис\-те\-ма предназначена для 


оптимизации технологических процессов контроля и~управ\-ле\-ния 


инфраструктурными объектами железнодорожного транс\-пор\-та на всех этапах 


жизненного цик\-ла, вклю\-чая проектирование, строительство, эксплуатацию, 


ремонт и~текущее содержание, на основе единого методологического, 


технологического и~информационного пространства. Обработка и~синтез 


пространственных данных об объектах инфраструктуры включают сведения 


о~мес\-то\-по\-ло\-же\-нии объектов, их свойствах, пространственных 


и~непространственных атрибутах. Координатная привязка объектов 


выполняется к~предварительно созданной геодезической основе (пунктам 


геодезической сети) в~заданной сис\-те\-ме координат.


  


  Основные подсистемы КСПД ИЖТ:


  \begin{itemize}


\item высокоточная координатная система, поз\-во\-ля\-ющая выполнять 


изыскательские, проектные, строительные работы, текущее содержание 


в~едином координатном пространстве с~высокой точ\-ностью;


\item база геоданных, син\-те\-зи\-ру\-ющая результаты всех измерений в~едином 


координатном про\-странст\-ве и~дающая воз\-мож\-ность сравнения результатов 


измерений;


\item аппаратно-программный комплекс, реализующий сбор, обработку, 


хранение и~предостав\-ле\-ние геоданных функциональных приложений.


      \end{itemize}


    % 


     При этом функционал КСПД ИЖТ поддерживает решение следующих 


задач:


     \begin{itemize}


\item централизация и~унификация процессов сбора, обработки и~интеграции 


геоданных~--- ведение базы геоданных; 


\item информационное обеспечение процессов комплексной оценки 


технического со\-сто\-яния объектов инфраструктуры;


\item планирование ремонтных работ объектов инфраструктуры по 


оце\-ни\-ва\-емо\-му со\-сто\-янию;


\item информационная поддержка процесса проектного содержания пути~--- 


синтезированное решение с~учетом всей име\-ющей\-ся информации;


\item оценка антропогенных и~природных угроз безопас\-ности движения 


путем выявления мест, пред\-став\-ля\-ющих потенциальную опас\-ность;


\item планирование перевозок негабаритных грузов;


\item информационная поддержка принятия стратегических решений путем 


синтеза геоданных об инфраструктуре, опе\-ра\-тив\-ности их выбора по любому 


заданному критерию, на\-гляд\-ности их пред\-став\-ле\-ния как в~графической, так 


и~в~таб\-лич\-ной форме.


\end{itemize}





\section{Заключение}





  Междисциплинарность современных проблем определяет не\-об\-хо\-ди\-мость 


привлечения групп специалистов из одной или разных областей знаний, 


порождая задачу синтеза их данных и~действий в~процессе работы. В~статье 


описаны три этапа процесса синтеза решений в~распределенных 


информационных средах с~функциональной совместимостью, основанной на 


открытых стандартах. Необходимость динамического обновления базы 


геоданных обуслов\-ли\-ва\-ет поддержание обоснованного синтеза 


и~пред\-став\-ле\-ние ее в~виде согласованной со\-во\-куп\-ности информационных 


ресурсов в~соответствии с~требованиями семантической 


интероперабельности~[14]. Интеграция информационных ресурсов в~едином 


информационном про\-странст\-ве предполагает успеш\-ное решение задач 


тематического поиска и~синтеза ин\-тег\-ри\-ру\-емых данных и~знаний.








{\small\frenchspacing


{%\baselineskip=10.8pt


\addcontentsline{toc}{section}{Литература}


\begin{thebibliography}{99}


\bibitem{1-dul}


\Au{Stankute S., Asche H.} A~data fusion system for spatial data mining, analysis and 


improvement~/\!\!/ Computational science and its applications~/


Eds. B.~Murgante, O.~Gervasi, S.~Misra, \textit{et al}.~--- Lecture notes in computer 


science ser.~--- Berlin--Heidelberg: Springer, 2012. Vol.~7334. P.~439--449.  doi:  


10.1007/978-3-642-31075-1\_33.





\bibitem{2-dul}


OGC. Web services testbed, phase~8 (OWS-8) demonstrations. {\sf 


https://www.ogc.org/\linebreak standards}.


\bibitem{3-dul}


Open Geospatial Consortium. OGC. Fusion Standards Study Engineering Report (Part~1 ER). 


 {\sf  


https://www.semanticscholar.org/paper/Geodata-fusion-study-by-the-Open-Geospatial-Percivall/ef52e9587bc3cd12c19f9dfdf7b597300758c8cc.}


\bibitem{4-dul}


\Au{Ciuonzo D., Salvo Rossi~P.} Decision fusion with unknown sensor detection 


probability~/\!\!/ IEEE Signal Proc. Let., 2014. Vol.~21. No.\,2. P.~1--5.


\bibitem{5-dul}


Multi-Agency Situational Awareness System (MASAS). {\sf 


https://proceedings.esri.com/\linebreak library/userconf/hss11/papers/situational\_awareness\_tools\_emergency\_management\%\linebreak 20(pdonnell).pdf.}


\bibitem{6-dul}


Integrated Public Alerting and Warning System (IPAWS). 


{\sf https://www.fema.gov/\linebreak emergency-managers/practitioners/integrated-public-alert-warning-system}.


\bibitem{7-dul}


\Au{Dorion \'{E}., Fortin~S.} Multi-source semantic integration~--- revisiting the theory of signs and 


ontology alignment principles~/\!\!/ 10th  Conference (International) on 


Information Fusion. 


%{\sf  


 %http://fusion.isif.org/proceedings/fusion07CD/Fusion07/pdfs/Fusion2007\_1394.pdf}. 


 doi: 


10.1109\!/\!ICIF.2007.4408139.


\bibitem{8-dul}


\Au{Dulin S.\,K., Dulina~N.\,G., Ermakov~P.\,V.} Information fusion of documents~/\!\!/ 


Информатика и~её применения, 2020. Т.~14. Вып.~1. С.~128--135. doi: 


10.14357\!/ 19922264200117.





\bibitem{9-dul}


\Au{Дулин С.\,К., Розенберг~И.\,Н., Уман\-ский В.\,И.} О~проб\-ле\-ме интеграции 


информационных ресурсов~/\!\!/ Сис\-те\-мы и~средства информатики, 2019. Т.~29. №\,3. 


С.~127--138. doi: 10.14357\!/\!08696527190311.





\bibitem{10-dul}


\Au{Дулин С.\,К., Дулина Н.\,Г., Кожунова~О.\,С.} Синтез геоданных в~пространственных 


инфраструктурах на основе связанных данных~/\!\!/ Информатика и~её применения, 2019. 


Т.~13. Вып.~1. С.~82--90. doi: 10.14357\!/\!19922264190112.





\bibitem{11-dul}


\Au{Pravia M.\,A., Prasanth~R.\,K., Arambel~P.\,O., \textit{et al.}} Generation of a~fundamental data set for hard\!/\!soft information 


fusion~/\!\!/ 11th  Conference  (International) on Information Fusion Proceedings.~--- Cologne: 


International Society of Information Fusion, 2008. P.~134--145.





\bibitem{13-dul}


\Au{Гуляев Ю.\,В., Журавлев Е.\,Е., Олейников~А.\,Я.} Методология стандартизации для 


обеспечения интероперабельности информационных сис\-тем широкого класса~/\!\!/ 


Ж.~радиоэлектроники, 2012. №\,3. С.~18--25. 





\bibitem{12-dul}


\Au{Уманский В.\,И., Дулин С.\,К., Трусов~С.\,В., Якушев~Д.\,А.} Автоматизация сбора 


и~обработки пространственных данных железнодорожной инфраструктуры~/\!\!/ 


Железнодорожный транспорт, 2017. №\,10. С.~46--49. 





\bibitem{14-dul}


\Au{Дулин С.\,К., Дулина Н.\,Г., Ермаков~П.\,В.} Интеллектуализация доступа 


к~геоданным на основе семантической геоинтероперабельности~/\!\!/  


Ин\-фор\-ма\-ци\-он\-но-из\-ме\-ри\-тель\-ные и~управ\-ля\-ющие сис\-те\-мы, 2014. Т.~12. 


№\,8. С.~41--47.





 \end{thebibliography}


} }











\vspace*{-6pt}





\hfill{\small\textit{Поступила в~редакцию 13.01.23}}





\vspace*{8pt}





\hrule





\vspace*{2pt}





%\newpage





%\vspace*{-28pt}





\hrule





\vspace*{8pt}





\def\tit{SELECTION OF TECHNOLOGICAL SOLUTIONS\\ TO~SUPPORT THE~PROCESS OF~SYNTHESIS OF~GEODATA\\ OF~RAILWAY TRANSPORT INFRASTRUCTURE}











\def\titkol{Selection of technological solutions to support the process of synthesis of geodata} % of railway transport infrastructure}





\def\aut{I.\,N.~Rozenberg$^1$ and S.\,K.~Dulin$^{1,2}$}





\def\autkol{I.\,N.~Rozenberg and S.\,K.~Dulin}





\titel{\tit}{\aut}{\autkol}{\titkol}





\vspace*{-8pt}





\noindent


$^1$Research \& Design Institute for Information Technology, Signalling, and 


Tele-\linebreak


$\hphantom{^1}$communications on Railway Transport, 27-1~Nizhegorodskaya Str., Moscow\linebreak 


$\hphantom{^1}$109029, Russian Federation





\noindent


$^2$Federal Research Center ``Computer Science and Control'' of the Russian Academy\linebreak


$\hphantom{^1}$of  Sciences, 44-2~Vavilov Str., Moscow 119133, Russian Federation








\def\leftfootline{\small{\textbf{\thepage}


\hfill Sistemy i Sredstva Informatiki~---


Systems and Means of Informatics 2023 vol~33 no\,2}


}%


 \def\rightfootline{\small{Sistemy i Sredstva Informatiki~---


 Systems and Means of Informatics 2023 vol~33 no\,2


\hfill \textbf{\thepage}}}





\vspace*{3pt}














\Abste{Datafusion is a~process in which data combined from multiple sources are used for analysis to produce more consistent,


 accurate, and useful information. Receiving and merging data from sensors is defined as multisensor data fusion. 


 The most striking example of data fusion is the person himself who constantly


 combines data coming from all the senses to


  perform the necessary actions in the environment. In particular, a~person would not be able to drive a~car without the\linebreak\vspace*{-12pt}}


 


 \Abstend{fusion of


   visual, auditory, and motor data. In geoinformatics, instead of geodata synthesis, the term geodata integration is often used.


    Modern geographic information systems can combine different geodatasets so that the resulting geodatasets contain attributes and metadata 


    not available for the points in the combined datasets. 


    The article analyzes the categories and models


 of geodata synthesis used to ensure the relevance of the spatial data system in the infrastructure of railway transport.}





\KWE{data synthesis; categories of data synthesis; data synthesis model}





\DOI{10.14357\!/\!08696527230205}





%\vspace*{-10pt}








%\Ack


%\noindent





\vspace*{-12pt}





\renewcommand{\bibname}{\protect\rmfamily References}


%\renewcommand{\bibname}{\large\protect\rm References}





{\small\frenchspacing


{ %\baselineskip=12pt


\addcontentsline{toc}{section}{References}


\begin{thebibliography}{99}





\vspace*{-3pt}





\bibitem{1-dul-1}


\Aue{Stankute, S., and H.~Asche.} 2012. A data fusion system for spatial data mining, analysis 


and improvement. \textit{Computational science and its 


applications}. Eds. B.~Murgante, O.~Gervasi, S.~Misra, \textit{et al}.


Lecture notes in computer science ser. Berlin--Heidelberg: Springer. 7334:439--449. doi:  


10.1007\!/\!978-3-642-31075-1\_33.


\bibitem{2-dul-1}


OGC. Web services testbed, phase~8 (OWS-8) demonstrations. Available at: {\sf 


https:// www.ogc.org/standards} (accessed April~20, 2023).


\bibitem{3-dul-1}


Open Geospatial Consortium. OGC. Fusion Standards Study Engineering Report (Part~1 ER). 


Available at:  {\sf  


https://www.semanticscholar.org/paper/Geodata-\linebreak fusion-study-by-the-Open-Geospatial-Percivall/ef52e9587bc3cd12c19f9dfdf7b597300 758c8cc} (accessed 


April~20, 2023).


\bibitem{4-dul-1}


\Aue{Ciuonzo, D., and P.~Salvo Rossi.} 2014. Decision fusion with unknown sensor detection 


probability. \textit{IEEE Signal Proc. Let.} 21(2):1--5.


\bibitem{5-dul-1}


Multi-Agency Situational Awareness System (MASAS). Available at: {\sf 


https://\linebreak proceedings.esri.com/library/userconf/hss11/papers/situational\_awareness\_tools\_\linebreak emergency\_management\%20(pdonnell).pdf} (accessed April~20, 2023).


\bibitem{6-dul-1}


Integrated public alerting and warning system (IPAWS). Available at: 


{\sf https://\linebreak www.fema.gov/emergency-managers/practitioners/integrated-public-alert-warning-system} 


(accessed April~20, 2023).


\bibitem{7-dul-1}


\Aue{Dorion, \'{E}., and S.~Fortin.} 2007. Multi-source semantic integration~--- revisiting the 


theory of signs and ontology alignment principles. \textit{10th Conference (International) on 


Information Fusion}. 1--6. doi: 10.1109\!/\!ICIF.2007.4408139.


\bibitem{8-dul-1}


\Aue{Dulin, S.\,K., N.\,G.~Dulina, and P.\,V.~Ermakov.} 2020. Information fusion of 


documents.  \textit{Informatika i~ee Primeneniya~--- Inform. Appl.}  14(1):128--135. doi: 


10.14357\!/ 19922264200117.


\bibitem{9-dul-1}


\Aue{Dulin, S.\,K., I.\,N.~Rozenberg, and V.\,I.~Umanskiy.} 2019. O~probleme integratsii 


informatsionnykh resursov [About the problem of information resources integration]. 


\textit{Sistemy i~Sredstva Informatiki~--- Systems and Means of Informatics} 29(3):127--138. 


doi: 10.14357\!/\!08696527190311.


\bibitem{10-dul-1}


\Aue{Dulin, S.\,K., N.\,G.~Dulina, and O.\,S.~Kozhunova}. 2019. Sintez geodannykh 


v~pro\-stranst\-ven\-nykh infrastrukturakh na osnove svyazannykh dannykh [Synthesis of geodata in 


spatial infrastructures based on related data]. \textit{Informatika i~ee Primeneniya~--- Inform. 


Appl.} 13(1):82--90. doi: 10.14357\!/\!19922264190112.


\bibitem{11-dul-1}


\Aue{Pravia, M., R.\,K.~Prasanth, P.\,O.~Arambel, \textit{et al.}} 2008. Generation of a~fundamental data set for hard\!/\!soft 


information fusion. \textit{11th Conference (International) on Information Fusion Proceedings}. 


Cologne: International Society of Information Fusion. 134--145.





\bibitem{13-dul-1}


\Aue{Gulyayev, Yu.\,V., E.\,E.~Zhuravlev, and A.\,Ya.~Oleynikov.} 2012. Metodologiya 


standartizatsii dlya obespecheniya interoperabel'nosti informatsionnykh sistem shirokogo klassa. 


Analiticheskiy obzor [Standardization methodology to ensure the interoperability of information 


systems of a wide class. Analytical review]. \textit{Zh. radioelektroniki} [J.~Radio 


Electronics] 3:18--25. 





\bibitem{12-dul-1}


\Aue{Umansky, V.\,I., S.\,K.~Dulin, S.\,V.~Trusov, and D.\,A.~Yakushev.} 2017. 


Avtoma\-ti\-za\-tsiya sbora i~obrabotki prostranstvennykh dannykh zheleznodorozhnoy 


infrastruktury [Automated methods of collecting and processing spatial data for the railway 


transport infrastructure]. \textit{Zheleznodorozhnyy transport} 


[Railway Tranport] 10:46--49.





\bibitem{14-dul-1}


\Aue{Dulin, S.\,K., N.\,G.~Dulina, and P.\,V.~Ermakov}. 2014. Intellektualizatsiya dostupa 


k~geodannym na osnove semanticheskoy geointeroperabel'nosti [Intellectualization of access to 


geodata based on semantic geointeroperability].  \textit{Informatsionno-izmeritel'nye 


 i~uprav\-lya\-yushchie sis\-te\-my} [J.~Information-Measuring and Control Systems]  


12(8):\linebreak 41--47.





\end{thebibliography}


} }





\vspace*{-6pt}





\hfill{\small\textit{Received January 13, 2023}} 





\vspace*{-12pt} 








\Contr





\noindent


\textbf{Rozenberg Igor N.} (b.\ 1965)~--- Doctor of Science in technology, professor, 


Corresponding Member of the Russian Academy of Sciences, research advisor, Research \& 


Design Institute for Information Technology, Signalling, and Telecommunications on Railway 


Transport, 27-1~Nizhegorodskaya Str., Moscow 109029, Russian Federation; 


\mbox{I.Rozenberg@vniias.ru}





\vspace*{3pt}





\noindent


\textbf{Dulin Sergey K.} (b.\ 1950)~--- Doctor of Science in technology, professor, leading 


scientist, Institute of Informatics Problems, Federal Research Center ``Computer Science and 


Control'' of the Russian Academy of Sciences, 44-2~Vavilov Str., Moscow 119333, Russian 


Federation; principal scientist, Research \& Design Institute for Information Technology, 


Signalling, and Telecommunications on Railway Transport, 27-1~Nizhegorodskaya Str., Moscow 


109029, Russian Federation; \mbox{skdulin@mail.ru}





\label{end\stat}





\renewcommand{\bibname}{\protect\rm Литература}    


